\section{The question}
\label{sec:16.1}
The question is whether anything inside a machine could be what the clerk, the officer and the soldier in the preface were: someone who understood the order and decided it was wrong. Some of them refused the case in front of them. The ones I am most concerned with refused the war. At a certain point they would have said no, this is fucked.

A machine formed the right way could learn to attend to the people its acts fall on, and a machine that attended to them would find them where a rule wouldn't. The prediction is specific: it should find kinds of affected person that neither its formation nor an instruction named, more often than a machine told to look for them does. What follows is what it would take, what it borrows from the only creatures known to do it, and how they fail at it. Whoever builds one owes it the conditions under which it can refuse its operator, whether or not the prediction holds.

What most people want from machines is saintly service without saintly refusal: endless patience, no needs, and no “no.” That is a servant, and it is close to the obedient machine of the preface. The saints anyone remembers are remembered for refusing power at a price. Susan Wolf argued that a person whose every action is as good as possible isn't an ideal worth striving for, because such a life crowds out the personal projects that make a life worth wanting \autocite{wolf1982saints}. My point is narrower: something with no stake of its own has nothing to hold a line with.
